\section{Related Work}
\label{sec:related}

\subsection{Coding Agent Taxonomy}
\label{sec:related:taxonomy}

AI coding tools can be organized by the degree of autonomous action they support (\Cref{tab:tool-taxonomy}). Inline completion tools such as GitHub Copilot \citep{chen2021evaluating} suggest code fragments within the editor without autonomous action. Chat-integrated products including Cursor and Windsurf add conversational interaction and multi-file edits but remain coupled to the IDE environment. Agentic CLI tools, including Claude Code, OpenAI's Codex CLI, and Aider \citep{gauthier2024aider}, operate from the command line and can autonomously execute shell commands, read and write files, and iterate on outputs within a single request. Fully autonomous systems like Devin, SWE-Agent \citep{yang2024sweagent}, and OpenHands \citep{wang2024openhands} aim for minimal human supervision, often in sandboxed cloud environments.

\begin{table}[t]
\centering
\caption{AI coding tool categories by degree of autonomous action.}
\label{tab:tool-taxonomy}
\small
\begin{tabularx}{\columnwidth}{>{\centering\arraybackslash}c >{\centering\arraybackslash}X >{\centering\arraybackslash}c}
\toprule
\textbf{Category} & \textbf{Examples} & \textbf{Pattern} \\
\midrule
Inline completion & Copilot, Tabnine & Editor plugin \\
Chat-integrated & Cursor, Windsurf, Cody & IDE-coupled product \\
Agentic CLI & Claude Code, Codex CLI, Aider & Tool-use loop \\
Fully autonomous & Devin, SWE-Agent, OpenHands & Sandbox + planning \\
\bottomrule
\end{tabularx}
\end{table}

Claude Code shares features with higher-autonomy agents (auto-mode classifier, background agent execution, remote environments) but retains interactive approval by default. Evaluation benchmarks such as SWE-Bench \citep{jimenez2024swebench} and HumanEval \citep{chen2021evaluating} have driven much of the academic focus on coding agents. This paper examines Claude Code's internal architecture from source code.

\subsection{Agent Architecture Patterns}
\label{sec:related:patterns}

Claude Code's core loop follows the ReAct pattern \citep{yao2022react}: the model generates reasoning and tool invocations, the harness executes actions, and results feed the next iteration. Toolformer \citep{schick2023toolformer} demonstrated that language models can learn to use tools. Claude Code uses up to 54 built-in tools and a layered permission system. The broader design space has been mapped by several surveys. \citet{weng2023agent} offered the now-standard decomposition into planning, memory, and tool use, and \citet{wang2024agentsurvey} catalogued early autonomous-agent work. \citet{xu2026agentsystems} frames the field around three recurring trade-offs (autonomy vs.\ controllability, latency vs.\ accuracy, capability vs.\ reliability) that recur throughout our analysis, and \citet{hu2025adas} casts agent design itself as a search problem over components, algorithms, and evaluation functions. This paper characterizes one specific point in that space.

A recent survey on code as agent harness frames code not only as model output but as the substrate for reasoning, action, environment modeling, and execution-based verification~\citep{ning2026codeharness}. That framing matches this paper's focus on the harness around the model rather than on benchmark scores alone.

Multi-agent orchestration frameworks such as AutoGen \citep{wu2024autogen}, LangChain, and CrewAI provide conversation-based agent coordination. Claude Code's subagent delegation (\Cref{sec:subagent}) includes permission override precedence, two-level permission scoping, and separate transcript files for each subagent. LATS \citep{zhou2024lats} unifies reasoning, acting, and planning in a tree-search framework. Claude Code's \code{plan} permission mode implements a simpler plan-then-execute approach.

Practitioner writing has converged on a handful of recurring patterns that Claude Code's architecture instantiates. Anthropic's own ``Building Effective Agents'' \citep{anthropic2024effective} distinguishes agents from workflows and argues for simple composable patterns over heavy frameworks. \citet{martin2026patterns} synthesizes seven patterns observed in production systems, including giving agents filesystem and shell access as a general-purpose action layer, and discovering actions on demand rather than loading every tool schema upfront. \citet{chase2025deepagents} observes that Claude Code's planning tool is ``basically a no-op'' whose value lies in keeping the agent on track rather than in performing any external computation. \citet{swyx2025agentengineering} argues that authority is the element academic frameworks most often leave out, calling trust ``the most overlooked element'' in production agent design, a gap the permission analysis in \Cref{sec:auth} attempts to close. \citet{huyen2025agents} makes the compound-error concern concrete: at 95\% per-step accuracy, a 100-step task succeeds only 0.6\% of the time, which motivates the per-step verification patterns we trace in \Cref{sec:turn} and \Cref{sec:auth}. OpenAI makes a similar point in its own writing on harness engineering, treating the design of the harness itself as a central engineering task~\citep{openai2026harness}. Recently, a new concept called \emph{loop engineering} has also emerged: instead of prompting the agent each turn, the developer builds a loop that finds work, assigns it, checks the result, and decides what comes next, assembled from the same primitives this paper documents, such as skills, MCP servers, and subagents~\citep{osmani2026loop}. It lets a single developer direct far more work autonomously and has drawn interest, including from Claude Code's own team, though reducing the person's turn-by-turn involvement also sharpens the long-term developer-capability concern that \Cref{sec:conclude} returns to.

\paragraph{Context management.}
\Cref{tab:context-comparison} presents a design-space taxonomy of context management approaches. Claude Code's five-layer compaction pipeline applies multiple strategies at different granularities before escalating, with cache-aware compression and virtual-view-on-read semantics. \citet{zhang2026ace} characterizes two failure modes that this design mitigates (summarization that drops domain details, and detail loss from iterative context rewriting), and instead proposes treating context as an ``evolving playbook'' that accumulates strategies over time. Claude Code's approach is consistent with that framing, since the CLAUDE.md hierarchy accumulates structured instructions rather than repeatedly summarizing them. \citet{hu2025memory} distinguishes context engineering from agent memory: context engineering handles transient assembly, while memory covers persistent factual knowledge and experiential traces. Claude Code's architecture separates the two in the same way, pairing a compaction pipeline with a file-based memory hierarchy.

\begin{table}[t]
\centering
\caption{Design space of context management approaches in LLM-based tools.}
\label{tab:context-comparison}
\small
\begin{tabularx}{\columnwidth}{>{\centering\arraybackslash}c >{\centering\arraybackslash}X >{\centering\arraybackslash}c}
\toprule
\textbf{Approach} & \textbf{Mechanism} & \textbf{Granularity} \\
\midrule
Simple truncation & Drop oldest messages & Coarse \\
Sliding window & Fixed-size recent history & Medium \\
RAG & Retrieve relevant snippets & Fine \\
Single summarization & One-pass compress & Coarse \\
Graduated compaction & Multi-layer pipeline & Very fine \\
\bottomrule
\end{tabularx}
\end{table}

\paragraph{Safety and permissions.}
Production coding agents adopt safety architectures that vary along three axes: \emph{approval model} (per-action prompting, classifier-mediated automation, or no prompting with post-hoc review), \emph{isolation boundary} (OS-level container, filesystem sandbox, permission-scoped tool pool, or none), and \emph{recovery mechanism} (version-control rollback, session-scoped permission reset, or checkpoint-based rewind). SWE-Agent and OpenHands~\citep{yang2024sweagent,wang2024openhands} rely primarily on Docker container isolation, providing environment-level sandboxing that constrains all agent actions. Codex CLI supports sandbox modes and approval policies for shell commands. Aider~\citep{gauthier2024aider} uses Git as its primary safety mechanism, making all changes reversible through version control. Claude Code combines per-action deny-first rules, an ML-based classifier for automated approval, optional shell sandboxing, and session-scoped permission non-restoration, layering multiple mechanisms rather than relying on a single isolation boundary.

Sandboxing reports make this boundary concrete. OpenAI's Windows sandbox discussion describes write access scoped to the working directory and configured writable roots, with explicit deny rules for selected read-only paths inside those roots~\citep{openai2026windowssandbox}. Cursor's cloud-agent report adds the cloud side of the same problem: network policy, secret redaction, and credential management become part of the agent environment~\citep{cursor2026cloudagents}. These details make the safety boundary an execution boundary, not only an approval prompt.

\paragraph{Protocols and extensibility.}
The Model Context Protocol that Claude Code uses as its primary external tool integration has become a de facto standard with a substantial ecosystem and a corresponding attack surface. \citet{hou2025mcpsurvey} catalogues thousands of community-developed MCP servers across 26 major directories and organizes MCP-specific threats into four attacker categories and sixteen scenarios, including tool poisoning, rug pulls, and cross-server shadowing. The permission and deny-rule machinery analyzed in \Cref{sec:auth} and the pre-filtering step in \Cref{sec:ext:assembly} can be read as the runtime side of the mitigations that survey calls for.

The same point applies to higher-level tool design. NSA guidance treats MCP security as a lifecycle issue that spans protocol design, runtime scheduling, integration with external services, and long-term monitoring~\citep{nsa2026mcpsecurity}. A proposal for agent-first tool design argues that agent-facing APIs should offer search, resolve, preview, execute, verify, and recover phases, not just human-oriented CRUD endpoints~\citep{pan2026agentfirstapi}. Read this way, MCP servers, plugins, skills, SDKs, and agent-to-agent protocols form a tool supply chain: they expand what an agent can do, but they also need identity, allowlists, versioning, audit logs, and revocation.

\paragraph{Software architecture.}
Layered architecture patterns \citep{garlan1993architecture} inform our five-layer decomposition. Role-based access control models \citep{sandhu1996rbac} provide theory for the permission mode system. Browser sandboxing \citep{reis2009isolating} is a similar per-process isolation approach. Multi-agent system theory \citep{wooldridge2009multiagent} helps explain subagent delegation.

\paragraph{Positioning.}
Prior work on coding agents has focused on benchmarks (how well agents solve tasks), frameworks (how to compose agents), and products (what users can do). This paper contributes a source-grounded design-space analysis of a production coding agent, using source-level analysis and architectural comparison to surface design choices and trade-offs. It draws on the software architecture case study tradition~\citep{garlan1993architecture} but applies it to an LLM-based agent by systematically identifying design questions, mapping alternatives, and contrasting Claude Code's choices with those of OpenClaw and Hermes Agent, two independent AI agent systems operating from different deployment contexts.
